\chapter{Case Study: the Intel 8086 Processor}

\begin{khulasa}[: what will we learn in this chapter?]
\begin{itemize}
    \item The 80x86 family, and the internal structure of the 8086.
    \item The 8086 registers, and segmented addressing (how to compute a physical address).
    \item Writing a first program, and how an instruction runs inside the processor.
    \item The instructions: transfer, arithmetic, logic, jumps, loops, then how an instruction becomes machine code.
    \item Input/output, the stack, procedures, and interrupts.
\end{itemize}
\end{khulasa}

%================================================
\section{The 80x86 family}
%================================================

In 1978 Intel released the 8086. It was a big success because IBM chose its brother, the 8088, for the first
personal computer (IBM PC) in 1981. Since then, every new Intel and AMD processor can still run 8086 programs.
This family is called \textbf{x86}.

\begin{tarif}
    \term{Backward compatibility}{compatibilité ascendante}: a new processor understands all the instructions
    of the older processors of its family. A program written for the 8086 in 1980 can run on a modern Core i7.
    This is why the x86 family has lived for more than 45 years.
\end{tarif}

\begin{center}
\small
\begin{tabular}{|c|c|c|c|l|}
    \hline
    \ras{Processor} & \ras{Year} & \ras{Bits} & \ras{Max memory} & \ras{What is new} \\ \hline
    8086 / 8088 & 1978 & 16 & 1 MB & first processor of the family \\ \hline
    80286 & 1982 & 16 & 16 MB & protected mode \\ \hline
    80386 & 1985 & 32 & 4 GB & first 32-bit processor (IA-32) \\ \hline
    80486 & 1989 & 32 & 4 GB & cache memory and floating-point unit inside \\ \hline
    Pentium & 1993 & 32 & 4 GB & runs two instructions at once \\ \hline
    Athlon 64 (AMD) & 2003 & 64 & $2^{48}$ in practice & the AMD64 (x86-64) architecture \\ \hline
    Core i7 and later & 2008 & 64 & huge & several cores on the same chip \\ \hline
\end{tabular}
\end{center}

\begin{tarif}
    \textbf{General features of the 8086}:
    \begin{itemize}
        \item a \textbf{16-bit} processor: its registers and its data bus are 16 bits wide.
        \item a \textbf{20-bit} address bus: it can address 1 MB of memory.
        \item about 29\,000 transistors, and a frequency between 5 and 10 MHz.
        \item about \textbf{100 different instructions} (a CISC processor).
    \end{itemize}
\end{tarif}

\begin{hal}
    The only difference between the 8086 and the 8088 is the \textbf{external data bus}: 16 bits in the 8086,
    only 8 bits in the 8088. Inside, they are identical. IBM chose the 8088 because 8-bit support chips were
    cheaper at that time.
\end{hal}

%================================================
\section{Inside the 8086}
%================================================

The 8086 is split into two units that work \textbf{at the same time}:

\begin{center}
\begin{tikzpicture}[r/.style={reg, minimum width=1.25cm, font=\codefont\small}]
    \draw[brandTeal, very thick, rounded corners=6pt, fill=brandTeal!8] (0.3,-3.2) rectangle (6.6,2.6);
    \node[font=\bfseries, text=tealDeep] at (3.45,2.25) {Execution Unit (EU)};
    \foreach \n/\y in {AX/1.4, BX/0.8, CX/0.2, DX/-0.4} \node[r] at (1.3,\y) {\n};
    \foreach \n/\y in {SI/1.4, DI/0.8, BP/0.2, SP/-0.4} \node[r] at (2.7,\y) {\n};
    \draw[brandGreen!70!black, very thick, fill=brandGreen!20] (4.1,1.2) -- (5.9,1.2) -- (5.5,-0.2) -- (4.5,-0.2) -- cycle;
    \node[font=\bfseries\small] at (5,0.5) {UAL};
    \node[r, minimum width=2.5cm, font=\small] at (5,-1.0) {FLAGS};
    \node[unit=brandOrange, font=\small, minimum width=2.6cm] at (2.2,-2.1) {control unit};
    \draw[brandOrange, very thick, rounded corners=6pt, fill=brandOrange!8] (-6.3,-3.2) rectangle (-0.3,2.6);
    \node[font=\bfseries, text=orangeDeep] at (-3.3,2.25) {Bus Interface Unit (BIU)};
    \node[r] at (-4.6,1.4) {CS}; \node[r] at (-4.6,0.8) {DS};
    \node[r] at (-4.6,0.2) {SS}; \node[r] at (-4.6,-0.4) {ES};
    \node[r] at (-4.6,-1.0) {IP};
    \node[draw=brandDark, fill=brandOrange!25, thick, minimum width=1.4cm, minimum height=0.6cm, font=\small] at (-2,1.1) {adder $\Sigma$};
    \node[font=\scriptsize, align=center] at (-2,0.35) {computes the\\physical address};
    \foreach \i in {0,...,5} \node[draw=brandDark, fill=white, minimum width=0.55cm, minimum height=0.5cm] at (-4.2+0.6*\i,-2.1) {};
    \node[font=\scriptsize] at (-2.7,-2.75) {instruction queue (6 bytes)};
    \draw[sahmT, <->] (-0.3,-2.1) -- (0.3,-2.1);
    \draw[brandDark, line width=4pt, {Stealth[length=4mm]}-{Stealth[length=4mm]}] (-6.3,0.2) -- (-7.8,0.2);
    \node[font=\small, align=center] at (-7.1,-0.5) {system\\bus};
\end{tikzpicture}
\end{center}

\begin{itemize}
    \item The \term{Bus Interface Unit}{unité d'interface de bus} (BIU): it talks with memory. It computes the
          addresses, fetches instructions in advance into a \textbf{queue} of 6 bytes, and reads and writes data.
    \item The \term{Execution Unit}{unité d'exécution} (EU): it takes the instructions from the queue, decodes
          them, and runs them with the arithmetic and logic unit.
\end{itemize}

\begin{tashbih}
    In a restaurant, the \textbf{waiter} (BIU) brings orders from the customers and puts them in a queue in front of
    the cook, while the \textbf{cook} (EU) prepares the current order. The cook never waits for the waiter, because
    the next order is always in the queue. This is the first idea of a pipeline!
\end{tashbih}

%================================================
\section{The 8086 registers}
%================================================

The 8086 has \textbf{14 registers} of 16 bits.

\subsection{General registers}
Four data registers. Each one can be split into two halves of 8 bits:

\begin{center}
\begin{tikzpicture}[h/.style={draw=brandDark, thick, minimum width=2.2cm, minimum height=0.7cm, font=\codefont\bfseries}]
    \foreach \r/\n/\d [count=\i from 0] in {A/AX/{accumulator: arithmetic, multiply, divide, input/output},
                                         B/BX/{base: keeps an address in memory},
                                         C/CX/{counter: number of loop turns},
                                         D/DX/{data: with AX for multiply and divide}} {
        \node[font=\codefont\bfseries, text=tealDeep] at (-1.6,-0.9*\i) {\n};
        \node[h, fill=brandTeal!25] (hi\i) at (0,-0.9*\i) {\r H};
        \node[h, fill=brandTeal!10] (lo\i) at (2.2,-0.9*\i) {\r L};
        \node[font=\small, anchor=west] at (3.6,-0.9*\i) {\d};
    }
    \node[font=\scriptsize] at (0,0.6) {bits 15 to 8};
    \node[font=\scriptsize] at (2.2,0.6) {bits 7 to 0};
\end{tikzpicture}
\end{center}

\begin{mithal}
    If AX = \hex{1234} then AH = \hex{12} and AL = \hex{34}. If we change AL to \hex{FF}, AX becomes \hex{12FF}.
\end{mithal}

Why split the registers? Because we often work with small values (a character, a number from 0 to 255).
Instead of wasting a whole 16-bit register for one character, we can put two characters in one register:
one in AH and one in AL.

\begin{intabih}
    These registers are ``general'' because we can use them in any calculation, but some instructions force a
    register: \code{loop} always uses CX, \code{mul} and \code{div} always use AX and DX, and the
    \code{int 21h} services expect the function number in AH. You must remember these roles.
\end{intabih}

\subsection{Pointer and index registers}
\begin{center}
\begin{tabular}{|c|l|l|}
    \hline
    \ras{Register} & \ras{Name} & \ras{Job} \\ \hline
    SI & Source Index & index of the source (arrays and strings) \\ \hline
    DI & Destination Index & index of the destination \\ \hline
    SP & Stack Pointer & points to the top of the stack \\ \hline
    BP & Base Pointer & to reach data inside the stack \\ \hline
\end{tabular}
\end{center}

\begin{mithal}
    To walk through an array of 10 elements, we put its start address in SI, then increase SI at each turn
    (\code{inc si}). At the first turn SI points to the first element, at the second turn to the second one,
    and so on. SI plays the role of the \textbf{finger} you put on the element you are reading now.
\end{mithal}

\subsection{Segment registers}
\begin{center}
\begin{tabular}{|c|l|l|}
    \hline
    \ras{Register} & \ras{Name} & \ras{Points to the start of\dots} \\ \hline
    CS & Code Segment & the program segment (the instructions) \\ \hline
    DS & Data Segment & the data segment (the variables) \\ \hline
    SS & Stack Segment & the stack segment \\ \hline
    ES & Extra Segment & an extra segment (for copying strings, for example) \\ \hline
\end{tabular}
\end{center}

\subsection{Instruction pointer and flags}
\begin{itemize}
    \item \textbf{IP} (Instruction Pointer): the address of the next instruction inside the code segment. It is the
          CO we saw in Chapter 2. We cannot change it directly with \code{mov}, only with jumps and calls.
    \item \textbf{FLAGS}: 9 bits are used out of 16:
\end{itemize}

\begin{center}
\begin{tikzpicture}[b/.style={draw=brandDark, minimum width=0.75cm, minimum height=0.7cm, font=\codefont\scriptsize\bfseries}]
    \foreach \n [count=\i from 0] in {{}, {}, {}, {}, OF, DF, IF, TF, SF, ZF, {}, AF, {}, PF, {}, CF} {
        \pgfmathtruncatemacro\bit{15-\i}
        \node[b, fill=brandTeal!15] at (0.75*\i, 0) {\n};
        \node[font=\codefont\tiny] at (0.75*\i, 0.55) {\bit};
    }
\end{tikzpicture}
\end{center}

\begin{center}
\small
\begin{tabularx}{\linewidth}{|c|R||c|R|}
    \hline
    \ras{Flag} & \ras{Meaning} & \ras{Flag} & \ras{Meaning} \\ \hline
    CF & carry or borrow out of the last bit & OF & overflow: wrong result for signed numbers \\ \hline
    PF & the number of 1s in the low byte is even & DF & direction for strings (0: forward) \\ \hline
    AF & carry from bit 3 to bit 4 & IF & hardware interrupts allowed \\ \hline
    ZF & the result is zero & TF & step-by-step execution (for debugging) \\ \hline
    SF & the result is negative (last bit is 1) & & \\ \hline
\end{tabularx}
\end{center}

%================================================
\section{Segmented addressing}
%================================================

\subsection{The problem}
The 8086 registers are 16 bits wide, so they can write addresses from 0 to $2^{16} - 1 = 65535$ only
(that is 64 KB). But the address bus is 20 bits wide (1 MB). How can we write a 20-bit address with 16-bit registers?

\subsection{The solution: segment and offset}
We use \textbf{two registers}: a \textbf{segment} register gives the start of a zone of memory, and the
\textbf{offset} gives the distance inside this zone. We write the logical address as \texttt{segment:offset}.

\begin{tarif}
    \centering
    $\text{physical address} = \text{segment} \times 16 + \text{offset}$
    \par\small
    Multiplying by 16 = adding one zero on the right in hexadecimal.
\end{tarif}

\begin{khatwa}[: the physical address for CS = 1000H and IP = 0234H]
    \begin{enumerate}
        \item Multiply the segment by 16 (add a zero): \hex{1000} becomes \hex{10000}.
        \item Add the offset:
        \begin{center}
        {\codefont
        \begin{tabular}{r}
              10000 \\
            + \ 0234 \\ \hline
              10234
        \end{tabular}}
        \end{center}
        \item The physical address is \hex{10234}, a 20-bit address (5 hex digits).
    \end{enumerate}
\end{khatwa}

\begin{center}
\begin{tikzpicture}
    \draw[thick, fill=brandLight] (0,0) rectangle (2.5,-7);
    \node[font=\small\bfseries] at (1.25,0.35) {memory (1 MB)};
    \node[font=\codefont\scriptsize, anchor=east] at (-0.1,-0.1) {00000};
    \node[font=\codefont\scriptsize, anchor=east] at (-0.1,-6.9) {FFFFF};
    \fill[brandTeal!30] (0,-2) rectangle (2.5,-4.5);
    \draw[thick] (0,-2) rectangle (2.5,-4.5);
    \node[font=\scriptsize, align=center] at (1.25,-3.4) {segment\\(64 KB)};
    \draw[sahmO] (4,-2) node[right, font=\small] {segment start: 10000} -- (2.5,-2);
    \draw[sahmT] (4,-2.8) node[right, font=\small] {address 10234} -- (2.5,-2.8);
    \draw[decorate, decoration={brace, amplitude=5pt, mirror}, thick] (-0.1,-2) -- (-0.1,-2.8)
        node[midway, left=6pt, font=\small] {offset 0234};
\end{tikzpicture}
\end{center}

\subsection{Why multiply by 16?}
If we used the segment as it is (without multiplying), we could not go beyond $\hex{FFFF} + \hex{FFFF} = \hex{1FFFE}$,
about 128 KB only. With the multiplication by 16, a segment can point anywhere in 1 MB: $\hex{FFFF} \times 16 = \hex{FFFF0}$.
So each segment always starts at an address that can be divided by 16. This block of 16 bytes is called a
\term{paragraph}{paragraphe}.

\begin{khatwa}[: which zone does the segment DS = 2000H cover?]
    \begin{enumerate}
        \item The smallest offset is \hex{0000}: $\hex{20000} + \hex{0000} = \hex{20000}$.
        \item The biggest offset is \hex{FFFF}: $\hex{20000} + \hex{FFFF} = \hex{2FFFF}$.
        \item The segment covers the zone from \hex{20000} to \hex{2FFFF}, that is $2^{16} = 65536$ bytes = 64 KB.
    \end{enumerate}
\end{khatwa}

\begin{tashbih}
    Imagine a town with many streets. The full address of a house is ``\textbf{street} number 12, \textbf{house}
    number 34''. The segment register is the street number (it gives the start of a zone), and the offset is the
    house number inside the street. The postman needs both pieces of information to find the house.
\end{tashbih}

\begin{intabih}
    One physical address can be written in several ways: 1000:0234 and 1023:0004 give the same address
    \hex{10234}. Check it yourself: $\hex{10230} + \hex{0004} = \hex{10234}$.
\end{intabih}

\subsection{Default pairs}
Each offset register is used by default with a given segment register:

\begin{center}
\begin{tabular}{|c|c|l|}
    \hline
    \ras{Offset} & \ras{Default segment} & \ras{Use} \\ \hline
    IP & CS & fetching instructions \\ \hline
    SP, BP & SS & the stack \\ \hline
    BX, SI, direct address & DS & variables \\ \hline
    DI (in string instructions) & ES & destination of a copy \\ \hline
\end{tabular}
\end{center}

\begin{hal}
    In the .COM programs we write in this chapter, the four segments CS = DS = SS = ES are equal: the whole
    program fits in one 64 KB segment, and starts at offset \hex{100}. This is the meaning of the line
    \code{org 100h} that we find in all our programs.
\end{hal}

%================================================
\section{8086 addressing modes}
%================================================

\begin{tarif}
    The \term{effective address}{adresse effective} (EA): the offset the processor computes for a memory operand.
    Then it computes the physical address $= \text{DS} \times 16 + \text{EA}$.
\end{tarif}

\begin{center}
\small
\begin{tabularx}{\linewidth}{|>{\centering\arraybackslash}p{2.8cm}|>{\centering\arraybackslash}p{3.4cm}|R|}
    \hline
    \ras{Mode} & \ras{Example} & \ras{Effective address EA} \\ \hline
    immediate & \code{mov ax, 5} & none: the value is in the instruction \\ \hline
    register & \code{mov ax, bx} & none: the value is in a register \\ \hline
    direct & \code{mov ax, [0200h]} & EA = \hex{0200} \\ \hline
    register indirect & \code{mov ax, [bx]} & EA = BX (or SI, DI, BP) \\ \hline
    based & \code{mov ax, [bx+4]} & EA = BX + 4 \\ \hline
    indexed & \code{mov ax, [si+4]} & EA = SI + 4 \\ \hline
    based indexed & \code{mov ax, [bx+si+4]} & EA = BX + SI + 4 \\ \hline
\end{tabularx}
\end{center}

\begin{khatwa}[: where does \code{mov ax, [bx+si+4]} read?]
    We suppose: DS = \hex{2000}, BX = \hex{0100}, SI = \hex{0020}.
    \begin{enumerate}
        \item Effective address: $\text{EA} = \hex{0100} + \hex{0020} + 4 = \hex{0124}$.
        \item Physical address: $\hex{2000} \times 16 + \hex{0124} = \hex{20000} + \hex{0124} = \hex{20124}$.
        \item The processor reads a word (two bytes): the byte at \hex{20124} goes into AL, and the byte at
              \hex{20125} goes into AH.
    \end{enumerate}
\end{khatwa}

\begin{khatwa}[: \code{mov ax, [bp+2]} uses the segment SS]
    We suppose SS = \hex{3000} and BP = \hex{FFF0}.
    \begin{enumerate}
        \item Effective address: $\text{EA} = \hex{FFF0} + 2 = \hex{FFF2}$.
        \item Because we used BP, the default segment is SS, not DS.
        \item Physical address: $\hex{30000} + \hex{FFF2} = \hex{3FFF2}$.
    \end{enumerate}
    This is important in procedures: BP is used to read the parameters that are in the stack.
\end{khatwa}

\begin{intabih}
    The 8086 stores a word in memory with the \textbf{small byte first} (little endian): the word \hex{1234}
    is stored as 34 then 12. Also, two memory operands in the same instruction are not allowed: \code{mov [a], [b]} is wrong.
\end{intabih}

%================================================
\section{A first program}
%================================================

\subsection{The tools}
We write the program in \textbf{EMU8086}: \emph{File $\rightarrow$ New $\rightarrow$ COM template}, then we press
\emph{Emulate} and run one instruction at a time with the \emph{Single step} button.

Every assembly program goes through three steps:
\begin{center}
\begin{tikzpicture}[node distance=1.9cm]
    \node[unit=brandTeal, minimum width=2.8cm] (src) {source file\\\code{hello.asm}};
    \node[unit=brandOrange, minimum width=2.8cm, right=of src] (bin) {executable file\\\code{hello.com}};
    \node[unit=brandGreen, minimum width=2.8cm, right=of bin] (run) {running\\in memory};
    \draw[sahmT] (src) -- node[above, font=\scriptsize, text=brandDark] {assembler} (bin);
    \draw[sahmT] (bin) -- node[above, font=\scriptsize, text=brandDark] {loading} (run);
\end{tikzpicture}
\end{center}
\begin{enumerate}
    \item \textbf{Writing}: we type the instructions in a text file (the source).
    \item \textbf{Assembling}: the assembler turns each instruction into machine-code bytes, and computes the
          addresses of the variables and labels.
    \item \textbf{Loading and running}: the system copies the bytes into memory starting at offset \hex{100},
          sets IP = \hex{100}, and the processor starts.
\end{enumerate}

\subsection{The program}
\asm{ch3/hello.asm}
\screen{ch3/hello.out}

\begin{sharh}
    \s{1}{A comment: everything after \code{;} is not run by the processor; it is an explanation for the reader.}
    \s{2}{\code{org 100h}: a .COM program starts at offset \hex{100}.}
    \s{4}{Put the address of the message \code{msg} into DX.}
    \s{5}{Choose DOS function \hex{09}: ``print the string whose address is in DX''.}
    \s{6}{\code{int 21h}: ask DOS to run the function chosen in AH.}
    \s{7}{\code{ret}: end of the program, back to the system.}
    \s{9}{Define the variable \code{msg} with \code{db} (define byte). The sign \code{\$} marks the end of the text.}
\end{sharh}

\begin{intabih}
    We put the data \textbf{after} \code{ret}. If we put it before the instructions, the processor would try
    to run the letters of the message as if they were instructions!
\end{intabih}

\begin{khatwa}[: the program in memory after loading]
    These are the real bytes made by the assembler:
    \begin{center}
    \small
    \begin{tabular}{|c|l|l|}
        \hline \ras{Offset} & \ras{Bytes} & \ras{Content} \\ \hline
        0100 & {\codefont 8D 16 09 01} & \code{lea dx, msg} (address \hex{0109}) \\ \hline
        0104 & {\codefont B4 09} & \code{mov ah, 09h} \\ \hline
        0106 & {\codefont CD 21} & \code{int 21h} \\ \hline
        0108 & {\codefont C3} & \code{ret} \\ \hline
        0109 & {\codefont 53 61 6C 61 6D \dots} & the letters \code{S a l a m \dots} \\ \hline
    \end{tabular}
    \end{center}
    The whole program is 9 bytes of instructions, and the message starts right after, at \hex{0109}. Notice
    that the assembler computed the address of \code{msg} by itself and put it inside the \code{lea} instruction.
\end{khatwa}

\subsection{Defining variables}
\begin{center}
\begin{tabular}{|l|l|}
    \hline
    \ras{Writing} & \ras{Meaning} \\ \hline
    \code{x db 5} & one byte with value 5 \\ \hline
    \code{y dw 1234h} & one word (two bytes) with value \hex{1234} \\ \hline
    \code{t db 14, 3, 27} & an array of 3 bytes \\ \hline
    \code{buf db 10 dup(0)} & 10 bytes, all 0 \\ \hline
    \code{s db 'Hi\$'} & a string of characters \\ \hline
\end{tabular}
\end{center}

%================================================
\section{How an instruction runs inside the 8086}
%================================================

Let us follow the instruction \code{mov cx, [si]} from the program \code{transfer.asm} (next section). In
memory, this instruction is written with the bytes 8B 0C at offset \hex{010C}, and SI = \hex{0112}.

\begin{center}
\small
\begin{tabularx}{\linewidth}{|>{\centering\arraybackslash}p{1.3cm}|>{\centering\arraybackslash}p{1.2cm}|R|}
    \hline
    \ras{Step} & \ras{Unit} & \ras{What happens} \\ \hline
    1 & BIU & computes the physical address $\text{CS} \times 16 + \text{IP}$, and fetches the bytes 8B 0C into the queue (often before we need them). \\ \hline
    2 & EU & takes 8B from the queue: ``move a word from memory to a register''. Then 0C: ``register CX, address in SI''. IP becomes \hex{010E}. \\ \hline
    3 & EU & computes the effective address: EA = SI = \hex{0112}. \\ \hline
    4 & BIU & computes $\text{DS} \times 16 + \hex{0112}$, sends it on the address bus, with a read signal. \\ \hline
    5 & BIU & the word \hex{ABCD} arrives on the data bus (CD then AB). \\ \hline
    6 & EU & writes the result into CX: CX = \hex{ABCD}. The flags do not change with \code{mov}. \\ \hline
\end{tabularx}
\end{center}

%================================================
\section{The instruction set in detail}
%================================================

We use these short names in the tables: reg = register, mem = memory box, imm = immediate value.

\subsection{Transfer instructions}
\begin{center}
\small
\begin{tabular}{|l|l|l|}
    \hline
    \ras{Instruction} & \ras{Meaning} & \ras{Example} \\ \hline
    \code{mov dest, src} & $\text{dest} \leftarrow \text{src}$ & \code{mov ax, bx} \\ \hline
    \code{xchg a, b} & swap the two values & \code{xchg ax, bx} \\ \hline
    \code{lea reg, mem} & reg gets the address of the variable & \code{lea si, tab} \\ \hline
    \code{push src} & put on the stack & \code{push ax} \\ \hline
    \code{pop dest} & take from the stack & \code{pop ax} \\ \hline
\end{tabular}
\end{center}

\asm{ch3/transfer.asm}

\begin{center}
\small
\begin{tabular}{|c|l|c|c|c|c|}
    \hline
    \ras{Line} & \ras{Instruction} & \ras{AX} & \ras{BX} & \ras{CX} & \ras{SI} \\ \hline
    4 & \code{mov ax, 1234h} & \cell{orangeDeep}{1234} & 0000 & 0000 & 0000 \\ \hline
    5 & \code{mov bl, ah} & 1234 & \cell{orangeDeep}{0012} & 0000 & 0000 \\ \hline
    6 & \code{mov bh, al} & 1234 & \cell{orangeDeep}{3412} & 0000 & 0000 \\ \hline
    7 & \code{xchg ax, bx} & \cell{orangeDeep}{3412} & \cell{orangeDeep}{1234} & 0000 & 0000 \\ \hline
    8 & \code{lea si, num} & 3412 & 1234 & 0000 & \cell{orangeDeep}{0112} \\ \hline
    9 & \code{mov cx, [si]} & 3412 & 1234 & \cell{orangeDeep}{ABCD} & 0112 \\ \hline
    10 & \code{mov [res], ax} & 3412 & 1234 & ABCD & 0112 \\ \hline
\end{tabular}
\end{center}
\begin{sharh}
    \s{5--6}{Copy the two halves of AX into BX, swapped: BH = 34 and BL = 12.}
    \s{8}{\code{lea} gives the \textbf{address} \hex{0112}, not the value.}
    \s{9}{Read the word stored at address \hex{0112}: \hex{ABCD}.}
    \s{10}{The only line that writes into memory: \code{res} becomes \hex{3412}.}
\end{sharh}

\subsection{Arithmetic instructions}
\begin{center}
\small
\begin{tabularx}{\linewidth}{|l|R|l|}
    \hline
    \ras{Instruction} & \ras{Meaning} & \ras{Flags} \\ \hline
    \code{add d, s} & $d \leftarrow d + s$ & all \\ \hline
    \code{adc d, s} & $d \leftarrow d + s + \text{CF}$ (to add big numbers) & all \\ \hline
    \code{sub d, s} & $d \leftarrow d - s$ & all \\ \hline
    \code{inc d} & $d \leftarrow d + 1$ & all except CF \\ \hline
    \code{dec d} & $d \leftarrow d - 1$ & all except CF \\ \hline
    \code{neg d} & $d \leftarrow -d$ (two's complement) & all \\ \hline
    \code{cmp a, b} & computes $a - b$ \textbf{without keeping the result} & all \\ \hline
    \code{mul s} & unsigned multiply (\code{imul} for signed) & CF, OF \\ \hline
    \code{div s} & unsigned divide (\code{idiv} for signed) & undefined \\ \hline
\end{tabularx}
\end{center}

\asm{ch3/arith.asm}

\begin{center}
\small
\begin{tabular}{|c|l|c|c|c|}
    \hline
    \ras{Line} & \ras{Instruction} & \ras{Result} & \ras{Decimal} & \ras{CF} \\ \hline
    4 & \code{mov al, 25} & AL = 19 & 25 & 0 \\ \hline
    5 & \code{add al, 17} & AL = 2A & 42 & 0 \\ \hline
    6 & \code{sub al, 2} & AL = 28 & 40 & 0 \\ \hline
    7 & \code{inc al} & AL = 29 & 41 & 0 \\ \hline
    8 & \code{dec al} & AL = 28 & 40 & 0 \\ \hline
    9 & \code{neg al} & AL = D8 & $-40$ & 1 \\ \hline
    10 & \code{mov bx, 1000} & BX = 03E8 & 1000 & 1 \\ \hline
    11 & \code{add bx, 2500} & BX = 0DAC & 3500 & 0 \\ \hline
\end{tabular}
\end{center}

\begin{khatwa}[: why is $-40$ written \hex{D8}? (two's complement)]
    \begin{enumerate}
        \item Write 40 in binary: \bin{0010\,1000}.
        \item Flip all the bits: \bin{1101\,0111}.
        \item Add 1: $\bin{1101\,1000} = \hex{D8}$.
    \end{enumerate}
    Check: $40 + (-40) = \hex{28} + \hex{D8} = \hex{100}$, and on 8 bits only 00 is left. Correct!
\end{khatwa}

\subsubsection{Multiplication and division}
The processor chooses the place depending on the size of the operand:
\begin{center}
\small
\begin{tabular}{|c|c|c|c|}
    \hline
    \ras{Instruction} & \ras{Operand} & \ras{Operation} & \ras{Result} \\ \hline
    \code{mul bl} & 8 bits & $\text{AL} \times \text{BL}$ & AX \\ \hline
    \code{mul cx} & 16 bits & $\text{AX} \times \text{CX}$ & DX:AX (32 bits) \\ \hline
    \code{div bl} & 8 bits & $\text{AX} \div \text{BL}$ & AL = quotient, AH = remainder \\ \hline
    \code{div cx} & 16 bits & $\text{DX:AX} \div \text{CX}$ & AX = quotient, DX = remainder \\ \hline
\end{tabular}
\end{center}

\asm{ch3/muldiv.asm}

\begin{sharh}
    \s{6}{$12 \times 11 = 132 = \hex{84}$: AX = \hex{0084}.}
    \s{10}{$300 \times 500 = 150000$, which is bigger than 65535, so it is split over two registers:
           DX = \hex{0002} and AX = \hex{49F0} (because $\hex{249F0} = 150000$).}
    \s{14}{$47 = 5 \times 9 + 2$: AL = 09 (quotient) and AH = 02 (remainder), so AX = \hex{0209}.}
\end{sharh}

\begin{intabih}
    Before \code{div cx} we must put DX = 0 if the number to divide is only in AX; otherwise the processor takes DX
    as part of the number. Dividing by 0 stops the program (interrupt number 0).
\end{intabih}

\subsubsection{Adding numbers bigger than 16 bits: \code{adc}}
The 8086 adds 16 bits at a time. To add two 32-bit numbers, we split each number into two words (high and
low). We add the low words with \code{add}, then the high words with \code{adc}, which adds the carry CF
from the first addition, exactly like we carry ``one'' to the next column when we add with a pen.

\asm{ch3/add32.asm}

\begin{center}
\small
\begin{tabular}{|c|l|c|c|c|}
    \hline
    \ras{Line} & \ras{Instruction} & \ras{DX} & \ras{AX} & \ras{CF} \\ \hline
    5--6 & load the first number & 0001 & FFFF & 0 \\ \hline
    7 & \code{add ax, 0001h} & 0001 & \cell{orangeDeep}{0000} & \cell{orangeDeep}{\textbf{1}} \\ \hline
    8 & \code{adc dx, 0000h} & \cell{orangeDeep}{0002} & 0000 & 0 \\ \hline
\end{tabular}
\end{center}
The result is DX:AX = 0002:0000 = \hex{00020000}. If we used \code{add} on line 8, the carry would be lost and
we would get \hex{00010000}, a wrong result.

\subsection{Logic instructions}
\begin{center}
\small
\begin{tabularx}{\linewidth}{|l|R|l|}
    \hline
    \ras{Instruction} & \ras{Usual use} & \ras{Example} \\ \hline
    \code{and} & switch bits off (a mask) & \code{and al, 0Fh} \\ \hline
    \code{or} & switch bits on & \code{or al, 80h} \\ \hline
    \code{xor} & flip bits, or clear a register & \code{xor cx, cx} \\ \hline
    \code{not} & flip all the bits & \code{not al} \\ \hline
    \code{test} & \code{and} without keeping the result (check a bit) & \code{test al, 1} \\ \hline
    \code{shl}, \code{shr} & shift: multiply or divide by 2 & \code{shl bl, 1} \\ \hline
    \code{rol}, \code{ror} & rotate: the bit that leaves comes back on the other side & \code{rol al, 1} \\ \hline
\end{tabularx}
\end{center}

\asm{ch3/logic.asm}

\begin{center}
\small
\begin{tabular}{|c|l|c|c|}
    \hline
    \ras{Line} & \ras{Instruction} & \ras{AL in binary} & \ras{Hex} \\ \hline
    4 & \code{mov al, 10110110b} & {\codefont 1011 0110} & B6 \\ \hline
    5 & \code{and al, 00001111b} & {\codefont 0000 0110} & 06 \\ \hline
    6 & \code{or  al, 10000000b} & {\codefont 1000 0110} & 86 \\ \hline
    7 & \code{xor al, 11111111b} & {\codefont 0111 1001} & 79 \\ \hline
    8 & \code{not al} & {\codefont 1000 0110} & 86 \\ \hline
\end{tabular}
\end{center}

\begin{sharh}
    \s{5}{\code{and} with 1 keeps the bit, and with 0 switches it off: we keep only the 4 low bits.}
    \s{6}{\code{or} with 1 switches the bit on: bit 7 is now 1.}
    \s{10--12}{BL: $5 \rightarrow 10 \rightarrow 20 \rightarrow 10$. A shift left multiplies by 2, and a shift right divides by 2.}
    \s{13}{\code{xor cx, cx}: every bit with itself gives 0, so CX = 0. Faster and shorter than \code{mov cx, 0}.}
\end{sharh}

\subsubsection{Practical uses of masks}
\begin{khatwa}[: is the number in AL even?]
    A number is even if bit 0 is 0. We write \code{test al, 1}: this instruction computes $\text{AL} \cdot 1$
    without changing AL.
    \begin{itemize}
        \item If AL = 6 = \bin{0000\,0110}: bit 0 is 0, the result is 0, and ZF = 1: the number is even.
        \item If AL = 7 = \bin{0000\,0111}: bit 0 is 1, the result is 1, and ZF = 0: the number is odd.
    \end{itemize}
    Then we write \code{jz even} to jump when the number is even.
\end{khatwa}

\begin{khatwa}[: turn a small letter into a capital letter with a mask]
    In the ASCII table: \code{'a'} $= \hex{61} = \bin{0110\,0001}$ and \code{'A'} $= \hex{41} = \bin{0100\,0001}$.
    The only difference is bit 5. So \code{and al, 11011111b} (that is \code{0DFh}) switches bit 5 off and turns
    any small letter into a capital letter, and leaves a capital letter as it is. This is another way than
    \code{sub al, 32}, which we see in the input section.
\end{khatwa}

\subsection{Flow control: jumps}
Usually the processor runs the instructions in order. A \textbf{jump} changes IP, so execution goes somewhere else.

\begin{center}
\small
\begin{tabularx}{\linewidth}{|l|R|l|}
    \hline
    \ras{Instruction} & \ras{Jumps if\dots} & \ras{Flags} \\ \hline
    \code{jmp L} & always (unconditional jump) & -- \\ \hline
    \code{je}, \code{jz} & equal / result is zero & ZF = 1 \\ \hline
    \code{jne}, \code{jnz} & not equal & ZF = 0 \\ \hline
    \code{jg}, \code{jge} & greater / greater or equal (\textbf{signed}) & ZF, SF, OF \\ \hline
    \code{jl}, \code{jle} & less / less or equal (\textbf{signed}) & SF, OF \\ \hline
    \code{ja}, \code{jae} & above / above or equal (\textbf{unsigned}) & CF, ZF \\ \hline
    \code{jb}, \code{jbe} & below / below or equal (\textbf{unsigned}) & CF, ZF \\ \hline
\end{tabularx}
\end{center}

\begin{tashbih}
    \code{cmp} is like \textbf{putting two things on a scale}, and a conditional jump is like saying: ``\textbf{if}
    the right side goes down, go right''. The scale (\code{cmp}) sets the flags, and the jump looks at them.
\end{tashbih}

\subsubsection{Example: if \dots\ else}
\asm{ch3/ifelse.asm}
\screen{ch3/ifelse.out}

\begin{center}
\begin{tikzpicture}[node distance=0.6cm]
    \node[unit] (c) {\code{cmp al, [b]}};
    \node[diamond, aspect=2.2, draw=brandOrange, fill=brandOrange!15, thick, below=of c, font=\small] (q) {$a > b$?};
    \node[unit=brandGreen, below right=0.4cm and 0.6cm of q] (y) {\code{lea dx, msg1}};
    \node[unit=brandPurple, below left=0.4cm and 0.6cm of q] (n) {\code{lea dx, msg2}};
    \node[unit, below=1.8cm of q] (p) {\code{int 21h} (print)};
    \draw[sahm] (c) -- (q);
    \draw[sahm] (q) -| node[pos=0.2, above, font=\small] {yes: \code{jg first}} (y);
    \draw[sahm] (q) -| node[pos=0.2, above, font=\small] {no} (n);
    \draw[sahm] (y) |- (p);
    \draw[sahm] (n) |- node[pos=0.7, below, font=\small] {\code{jmp show}} (p);
\end{tikzpicture}
\end{center}

\begin{sharh}
    \s{5}{\code{cmp} computes $7 - 12 = -5$ and does not keep it, but it sets the flags (SF = 1).}
    \s{6}{\code{jg}: is $a > b$? No, so we do not jump and continue at line 7.}
    \s{7--8}{Choose the second message, then \code{jmp show} is \textbf{needed} so that line 10 is not run.}
    \s{12--13}{The printing is shared by the two cases.}
\end{sharh}

\begin{intabih}
    We often forget \code{jmp show} on line 8! Without it, the program also runs \code{lea dx, msg1}, so it
    always prints the first message.
\end{intabih}

\subsubsection{Signed or unsigned? \code{jg} or \code{ja}}
The same bits can be read in two ways (we saw this in Chapter 2). AL = \hex{C8} is 200 unsigned, and $-56$
signed. So the processor has two kinds of jumps: \code{ja}/\code{jb} for unsigned numbers, and \code{jg}/\code{jl}
for signed numbers.

\asm{ch3/signed.asm}
\screen{ch3/signed.out}

\begin{sharh}
    \s{4--6}{\code{ja}: compare 200 with 100 as positive numbers. $200 > 100$, so we jump.}
    \s{14--16}{\code{jg}: the same bits \hex{C8}, but read as signed: $-56 > 100$? No, so we do not jump.}
\end{sharh}

\begin{intabih}
    Choose the jump from the \textbf{kind of data}: ages, addresses, character codes, sizes are unsigned
    (\code{ja}, \code{jb}). Temperatures, balances, differences are signed (\code{jg}, \code{jl}).
\end{intabih}

\subsubsection{From the algorithm to assembly}
Each control structure from the algorithms course has a fixed ``translation'' in assembly. The idea is
always the same: we \textbf{compare} with \code{cmp}, then we \textbf{jump} depending on the result.

\newcommand{\stk}[1]{\begin{tabular}[t]{@{}l@{}}#1\end{tabular}}
\begin{center}
\small
\begin{tabular}{|l|l|l|}
    \hline
    \ras{Structure} & \ras{In C} & \ras{In assembly} \\ \hline
    if & \code{if (a == b) x = 1;} &
        \stk{\code{    cmp a, b}\\ \code{    jne end}\\ \code{    mov x, 1}\\ \code{end:}} \\ \hline
    while & \code{while (cx > 0) \{ ... \}} &
        \stk{\code{L:  cmp cx, 0}\\ \code{    jle end}\\ \code{    ...}\\ \code{    jmp L}\\ \code{end:}} \\ \hline
    for & \code{for (cx = 10; cx > 0; cx--)} &
        \stk{\code{    mov cx, 10}\\ \code{L:  ...}\\ \code{    loop L}} \\ \hline
\end{tabular}
\end{center}

\begin{intabih}
    In the translation of \code{if}, we jump with the \textbf{opposite} condition (\code{jne} instead of \code{je}).
    We tell the processor: ``if they are not equal, skip the assignment''. This makes the program shorter.
\end{intabih}

\subsection{Loops and repetition}
\begin{tarif}
    \code{loop L} is like two instructions: \code{dec cx}, then ``if CX $\neq 0$, jump to L''.
    So CX is the loop counter.
\end{tarif}

\subsubsection{The sum of the numbers from 1 to 10}
\asm{ch3/loopsum.asm}

\begin{center}
\small
\begin{tabular}{|c|c|c|c|c|c|c|c|c|c|c|}
    \hline
    \ras{Turn} & 1 & 2 & 3 & 4 & 5 & 6 & 7 & 8 & 9 & 10 \\ \hline
    \ras{CX before \code{add}} & 10 & 9 & 8 & 7 & 6 & 5 & 4 & 3 & 2 & 1 \\ \hline
    \ras{AX after \code{add}} & 10 & 19 & 27 & 34 & 40 & 45 & 49 & 52 & 54 & \cell{orangeDeep}{\textbf{55}} \\ \hline
\end{tabular}
\end{center}
After the tenth turn CX = 0, so the loop stops, and AX = \hex{0037} = 55. Notice that we add backwards:
$10 + 9 + \dots + 1$, and the result is the same.

\begin{intabih}
    If we enter the loop with CX = 0, \code{loop} first decreases CX, which becomes \hex{FFFF}, and then repeats
    65535 more times! To avoid this, we put \code{jcxz end} before the loop: ``if CX = 0, jump to the end''.
\end{intabih}

\subsubsection{Finding the biggest value in an array}
\asm{ch3/maxarray.asm}

\begin{center}
\small
\begin{tabular}{|c|c|c|c|c|}
    \hline
    \ras{SI points to} & \ras{Element} & \ras{Element $>$ AL?} & \ras{AL after} & \ras{CX after \code{loop}} \\ \hline
    tab[0] & 14 & no (equal) & 14 & 5 \\ \hline
    tab[1] & 3 & no & 14 & 4 \\ \hline
    tab[2] & 27 & yes & 27 & 3 \\ \hline
    tab[3] & 9 & no & 27 & 2 \\ \hline
    tab[4] & 31 & yes & \cell{orangeDeep}{\textbf{31}} & 1 \\ \hline
    tab[5] & 8 & no & 31 & 0 \\ \hline
\end{tabular}
\end{center}
The result: AL = \hex{1F} = 31, stored in \code{max}. We used \code{jbe} because the values are \textbf{unsigned}.

\subsubsection{Repetition with the \code{rep} prefix}
String instructions work on DS:SI (the source) and ES:DI (the destination), and the prefix \code{rep}
repeats them CX times.

\asm{ch3/repmovs.asm}
\screen{ch3/repmovs.out}

\begin{sharh}
    \s{7}{\code{cld}: set DF = 0, so SI and DI increase after each copy.}
    \s{8}{\code{rep movsb} copies one byte from [SI] to [DI], and repeats 6 times: the five letters and \code{\$}.}
\end{sharh}

\subsection{Generating machine code}
The assembler turns each instruction into bytes. These are the real bytes of the program \code{machine.asm}:

\asm{ch3/machine.asm}

\begin{center}
\small
\begin{tabular}{|c|l|l|}
    \hline
    \ras{Address} & \ras{Machine code (hex)} & \ras{Instruction} \\ \hline
    0100 & {\codefont B8 34 12} & \code{mov ax, 1234h} \\ \hline
    0103 & {\codefont 89 C3} & \code{mov bx, ax} \\ \hline
    0105 & {\codefont 04 05} & \code{add al, 5} \\ \hline
    0107 & {\codefont 88 07} & \code{mov [bx], al} \\ \hline
    0109 & {\codefont 41} & \code{inc cx} \\ \hline
    010A & {\codefont EB 01} & \code{jmp fin} \\ \hline
    010C & {\codefont 90} & \code{nop} \\ \hline
    010D & {\codefont C3} & \code{ret} \\ \hline
\end{tabular}
\end{center}

\begin{khatwa}[: decoding B8 34 12]
    \begin{enumerate}
        \item B8 = \bin{1011\,1000}: the first five bits 10111 mean ``\code{mov} an immediate value into a 16-bit register'',
              and the last three bits 000 are the number of the register AX.
        \item 34 12: the value \hex{1234}, written with the small byte first (little endian).
    \end{enumerate}
\end{khatwa}

\begin{khatwa}[: decoding 89 C3 (\code{mov bx, ax})]
    The first byte 89 = \bin{100010\,0\,1}:
    \begin{itemize}
        \item 100010: the opcode of \code{mov} between a register and a register or memory.
        \item d = 0: the register named in the \emph{reg} field is the \textbf{source}.
        \item w = 1: we work on 16 bits (if it were 0 we would work on 8 bits).
    \end{itemize}
    The second byte C3 = \bin{11\,000\,011} is called ModR/M:
    \begin{center}
    \begin{tikzpicture}[f/.style={draw=brandDark, thick, minimum height=0.7cm, font=\codefont\bfseries}]
        \node[f, minimum width=1.6cm, fill=brandOrange!20] (m) at (0,0) {11};
        \node[f, minimum width=2.2cm, fill=brandTeal!20, anchor=west] (r) at (m.east) {000};
        \node[f, minimum width=2.2cm, fill=brandGreen!20, anchor=west] (rm) at (r.east) {011};
        \node[below=0.1cm of m, font=\small] {mod: register};
        \node[below=0.1cm of r, font=\small] {reg = AX};
        \node[below=0.1cm of rm, font=\small] {r/m = BX};
    \end{tikzpicture}
    \end{center}
    So: move AX (the source) into BX.
\end{khatwa}

\begin{center}
\small
\begin{tabular}{|c|c|c|c|c|c|c|c|c|}
    \hline
    \ras{Number} & 000 & 001 & 010 & 011 & 100 & 101 & 110 & 111 \\ \hline
    \ras{w = 1} & AX & CX & DX & BX & SP & BP & SI & DI \\ \hline
    \ras{w = 0} & AL & CL & DL & BL & AH & CH & DH & BH \\ \hline
\end{tabular}
\end{center}

\begin{mithal}[: 88 07 and EB 01]
    88 = \bin{100010\,0\,0}: \code{mov} on 8 bits. 07 = \bin{00\,000\,111}: mod = 00 (memory without offset),
    reg = 000 (AL), r/m = 111 ([BX]). So \code{mov [bx], al}.\\[3pt]
    EB 01: a short jump, and the offset is \textbf{relative}: after reading \code{jmp}, IP = \hex{010C}, and we add 1
    to reach \hex{010D} (\code{fin}). This is how we jumped over \code{nop}.
\end{mithal}

%================================================
\section{Input and output}
%================================================

\subsection{Ports: \code{in} and \code{out}}
The processor talks with devices through numbered \term{ports}{ports d'entrée/sortie}, from 0 to \hex{FFFF}.
\begin{itemize}
    \item \code{in al, 60h}: read one byte from port \hex{60} (the keyboard in a PC).
    \item \code{out 61h, al}: send AL to port \hex{61} (the speaker).
\end{itemize}
Normal programs rarely use ports directly; they ask the operating system through interrupts.

\subsection{DOS services: the \code{int 21h} interrupt}
\begin{center}
\small
\begin{tabularx}{\linewidth}{|c|R|R|}
    \hline
    \ras{AH} & \ras{Function} & \ras{Inputs and outputs} \\ \hline
    \hex{01} & read a key from the keyboard and show it & output: the character in AL \\ \hline
    \hex{02} & print a character & input: the character in DL \\ \hline
    \hex{09} & print a string ending with \code{\$} & input: its address in DX \\ \hline
    \hex{4C} & end the program & AL = exit code \\ \hline
\end{tabularx}
\end{center}

\asm{ch3/io.asm}
\screen{ch3/io.out}

\begin{sharh}
    \s{7--8}{Wait for the user to press a key (here \code{m}); its ASCII code arrives in AL: 109.}
    \s{9}{In the ASCII table, a capital letter is 32 less than the small one: $109 - 32 = 77$, which is the code of \code{'M'}.}
    \s{10}{Keep the letter in BL, because the printing on line 13 changes AL.}
    \s{14--16}{Function \hex{02} prints the character that is in DL.}
    \s{20}{\code{13, 10}: back to the start of the line, then a new line.}
\end{sharh}

%================================================
\section{The stack}
%================================================

\begin{tarif}
    The \term{stack}{pile}: a zone of memory where we put values and take them back with the rule
    \textbf{last in, first out} (LIFO). Its top is given by SS:SP.
\end{tarif}

\begin{tashbih}
    The stack is like a \textbf{pile of plates}: we put the new plate on top of the pile, and we always take the plate on top.
\end{tashbih}

\begin{itemize}
    \item \code{push ax}: \nql{SP}{SP - 2}, then write AX at [SP].
    \item \code{pop ax}: read [SP] into AX, then \nql{SP}{SP + 2}.
\end{itemize}

\begin{intabih}
    In the 8086 the stack \textbf{grows towards small addresses}: each \code{push} decreases SP. And we always
    put and take \textbf{words} (16 bits).
\end{intabih}

\asm{ch3/stack.asm}

\begin{center}
\begin{tikzpicture}[c/.style={draw=brandDark, fill=white, minimum width=1.6cm, minimum height=0.55cm, font=\codefont\small},
    a/.style={font=\codefont\scriptsize, text=tealDeep}]
    \foreach \t/\x/\v/\w/\sp in {{start}/0/{}/{}/FFFE, {after \code{push ax}}/3.3/1111/{}/FFFC,
                                {after \code{push bx}}/6.6/1111/2222/FFFA, {after \code{pop ax}}/9.9/1111/{}/FFFC} {
        \node[font=\small\bfseries] at (\x,1.1) {\t};
        \node[c] (u) at (\x,0.3) {\v};
        \node[c] (d) at (\x,-0.25) {\w};
        \node[a, left=0.05cm of u] {FFFC};
        \node[a, left=0.05cm of d] {FFFA};
        \node[font=\codefont\small] at (\x,-1) {SP = \sp};
    }
\end{tikzpicture}
\end{center}

\begin{center}
\small
\begin{tabular}{|c|l|c|c|c|}
    \hline
    \ras{Line} & \ras{Instruction} & \ras{AX} & \ras{BX} & \ras{SP} \\ \hline
    4--5 & \code{mov ax}, \code{mov bx} & 1111 & 2222 & FFFE \\ \hline
    6 & \code{push ax} & 1111 & 2222 & \cell{orangeDeep}{FFFC} \\ \hline
    7 & \code{push bx} & 1111 & 2222 & \cell{orangeDeep}{FFFA} \\ \hline
    8 & \code{pop ax} & \cell{orangeDeep}{2222} & 2222 & FFFC \\ \hline
    9 & \code{pop bx} & 2222 & \cell{orangeDeep}{1111} & FFFE \\ \hline
\end{tabular}
\end{center}
Because the last value in (\hex{2222}) is the first value out, AX and BX have swapped their values.

\subsection{Why do we need a stack?}
\begin{enumerate}
    \item \textbf{To keep return addresses}: when we call a procedure, the processor keeps in the stack where it must come back.
    \item \textbf{To keep registers for a while}: a procedure needs CX, but CX is used by a loop? We write
          \code{push cx} at its start and \code{pop cx} at its end, so the loop counter is not lost.
    \item \textbf{To pass parameters} to procedures (we see it soon).
    \item \textbf{Interrupts}: the processor keeps its state there before going to the interrupt handler.
\end{enumerate}

\begin{intabih}
    The number of \code{pop} must equal the number of \code{push}, in the opposite order. \code{push ax} then
    \code{push bx} must be matched by \code{pop bx} then \code{pop ax}. If you forget one \code{pop} inside a
    procedure, \code{ret} takes a wrong value as the return address, and the program jumps to a random place.
\end{intabih}

%================================================
\section{Functions and procedures}
%================================================

\begin{tarif}
    A \term{procedure}{procédure, sous-programme}: a group of instructions with a name. We call it with
    \code{call} and come back from it with \code{ret}. We write it once and use it many times.
\end{tarif}

\subsection{How does \code{ret} know where to go back?}
\begin{itemize}
    \item \code{call P}: puts on the stack \textbf{the address of the instruction after it} (IP), then jumps to P.
    \item \code{ret}: takes this address from the stack and puts it into IP, so execution goes back after \code{call}.
\end{itemize}

\begin{khatwa}[: the two calls in the next program (real addresses)]
    \begin{enumerate}
        \item The first \code{call print\_num} is at offset \hex{0103} and is 3 bytes long, so the next instruction is
              at \hex{0106}. The processor puts \hex{0106} on the stack, then jumps to \code{print\_num}.
        \item At the end of the procedure, \code{ret} takes \hex{0106} from the stack: execution goes back to \code{mov ax, 7}.
        \item The second call is at \hex{0109}, so it puts \hex{010C}, and this time \code{ret} goes back to \hex{010C}.
    \end{enumerate}
    The same procedure goes back to two different places, because the return address is kept in the stack each time.
\end{khatwa}

\subsection{A procedure that prints a number in decimal}
\code{int 21h} only prints characters. To print the number 2026, we must take out its digits one by one:

\begin{khatwa}[: taking out the digits of 2026]
    \begin{center}
    \begin{tabular}{|c|c|c|c|}
        \hline
        \ras{Operation} & \ras{Quotient (AX)} & \ras{Remainder (DX)} & \ras{Stack} \\ \hline
        $2026 \div 10$ & 202 & 6 & 6 \\ \hline
        $202 \div 10$ & 20 & 2 & 6, 2 \\ \hline
        $20 \div 10$ & 2 & 0 & 6, 2, 0 \\ \hline
        $2 \div 10$ & 0 (we stop) & 2 & 6, 2, 0, 2 \\ \hline
    \end{tabular}
    \end{center}
    The digits come out backwards (6, 2, 0, 2). We put them on the stack, then take them back: they come out in
    the right order 2, 0, 2, 6.
\end{khatwa}

\asm{ch3/procnum.asm}
\screen{ch3/procnum.out}

\begin{sharh}
    \s{4--5}{Put the number in AX (the parameter), then call the procedure.}
    \s{10}{\code{proc} and \code{endp} mark the start and the end of the procedure.}
    \s{14}{DX = 0 is needed before \code{div bx} (a 16-bit division divides DX:AX).}
    \s{16--17}{Put the digit on the stack, and count the digits in CX.}
    \s{21--25}{CX times: take a digit, turn it into a character by adding \code{'0'} (48), and print it.}
    \s{28}{\code{ret} goes back to line 6 the first time, and to line 8 the second time.}
\end{sharh}

\subsection{Passing parameters on the stack}
\asm{ch3/procstack.asm}

After \code{mov bp, sp} on line 14, the stack looks like this (values from the real run):

\begin{center}
\begin{tikzpicture}[c/.style={draw=brandDark, fill=white, minimum width=3.6cm, minimum height=0.6cm, font=\small},
    a/.style={font=\codefont\small, text=tealDeep}]
    \foreach \ad/\v/\l [count=\i from 0] in {FFFC/{30 (2nd parameter)}/{BP+6}, FFFA/{12 (1st parameter)}/{BP+4},
                                           FFF8/{return address 010B}/{BP+2}, FFF6/{old BP}/{BP}} {
        \node[c] (n\i) at (0,-0.6*\i) {\v};
        \node[a, left=0.15cm of n\i] {\ad};
        \node[font=\codefont\small\bfseries, right=0.15cm of n\i] {\l};
    }
    \draw[sahmO] (4.6,-1.8) node[right, font=\small] {SP = BP = FFF6} -- (3.0,-1.8);
\end{tikzpicture}
\end{center}

\begin{sharh}
    \s{4--7}{Put the two parameters on the stack: SP goes from \hex{FFFE} to \hex{FFFA}.}
    \s{8}{\code{call} puts the return address \hex{010B}: SP = \hex{FFF8}.}
    \s{13--14}{Save BP, then make BP stay fixed on the top of the stack (\hex{FFF6}).}
    \s{15--16}{The first parameter is at [BP+4], the second at [BP+6]: $12 + 30 = 42$.}
    \s{18}{\code{ret 4}: go back, then remove 4 bytes (the two parameters) from the stack. SP goes back to \hex{FFFE}.}
\end{sharh}

%================================================
\section{Interrupts}
%================================================

\begin{tarif}
    An \term{interrupt}{interruption}: an event that stops the current program for a moment, so the processor
    runs a special program (the \textbf{interrupt handler}), then comes back to continue where it stopped.
\end{tarif}

\begin{tashbih}
    You are reading a book and the \textbf{doorbell} rings. You put a bookmark in the page, open the door, then
    come back and continue reading from the same page. The bell is the interrupt, and the bookmark is what the
    processor keeps in the stack.
\end{tashbih}

\subsection{Kinds of interrupts}
\begin{center}
\small
\begin{tabularx}{\linewidth}{|>{\centering\arraybackslash}p{3cm}|R|R|}
    \hline
    \ras{Kind} & \ras{Who causes it?} & \ras{Examples} \\ \hline
    hardware & an outside device sends a signal on the INTR or NMI pin & a key press, the clock, the printer is ready \\ \hline
    software & the instruction \code{int n} in the program & \code{int 21h} for DOS services \\ \hline
    exceptions & the processor itself, after an error & division by zero (interrupt 0) \\ \hline
\end{tabularx}
\end{center}

\subsection{The interrupt vector table}
At the start of memory (from \hex{00000} to \hex{003FF}) there is a table of 256 entries. Each interrupt has
4 bytes: the offset, then the segment, of its handler.

\begin{tarif}
    \centering
    The vector of interrupt number $n$ is at address: \quad $4 \times n$
\end{tarif}

\begin{mithal}
    The vector of \code{int 21h}: $\hex{21} \times 4 = 33 \times 4 = 132 = \hex{84}$. So the address of the DOS
    handler is kept from \hex{00084} to \hex{00087}.
\end{mithal}

\subsection{What does the processor do on \code{int n}?}
\begin{enumerate}
    \item It puts FLAGS on the stack, then switches IF and TF off.
    \item It puts CS, then IP, on the stack (the bookmark in the book).
    \item It reads the new IP and CS from the entry $n \times 4$, so it jumps to the interrupt handler.
    \item At the end of the handler, \code{iret} takes IP, CS and FLAGS back from the stack, and the program continues.
\end{enumerate}

\subsection{A hardware interrupt, step by step}
\begin{khatwa}[: you press a key while a program runs]
    \begin{enumerate}
        \item The keyboard sends a signal to the \term{interrupt controller}{contrôleur d'interruptions} (PIC 8259).
        \item The controller sends a signal on the INTR input of the processor.
        \item If IF = 1, the processor first finishes the \textbf{current instruction}, then accepts the interrupt.
        \item The controller gives the processor the interrupt number (\hex{09} for the keyboard).
        \item The processor saves FLAGS, CS and IP, then jumps to the address written at $4 \times 9 = \hex{24}$.
        \item The handler reads the key from port \hex{60}, keeps it, then runs \code{iret}.
        \item The program continues from where it stopped, without ``noticing'' anything.
    \end{enumerate}
\end{khatwa}

\begin{tarif}
    \begin{itemize}
        \item \code{cli} (clear IF) blocks hardware interrupts for a while, and \code{sti} (set IF) allows them again.
              We use them when a sensitive part of the program must not be interrupted.
        \item The NMI input (Non Maskable Interrupt): an interrupt that \textbf{cannot be blocked}, used for serious
              errors such as a memory fault or a power failure.
    \end{itemize}
\end{tarif}

\subsection{A program: our own interrupt}
\asm{ch3/myint.asm}
\screen{ch3/myint.out}

\begin{sharh}
    \s{4--5}{ES = 0: we want to write in the vector table, which is in segment 0.}
    \s{6--8}{Write the address of \code{handler} in entry $\hex{60} \times 4 = \hex{180}$: the offset, then the segment.}
    \s{10--11}{Each \code{int 60h} runs \code{handler}: the first time it prints \code{A}, the second time \code{B}.}
    \s{19}{\code{iret} and not \code{ret}, because the interrupt also put FLAGS on the stack.}
\end{sharh}

%================================================
\section{Exercises}
%================================================

\begin{tamrin}
    Compute the physical address in each case: (a) DS = 3000H, BX = 0050H with \code{mov al, [bx]}.
    (b) SS = 4000H, SP = FFFEH.
\end{tamrin}
\begin{hall}
    (a) $\hex{30000} + \hex{0050} = \hex{30050}$. \quad (b) $\hex{40000} + \hex{FFFE} = \hex{4FFFE}$.
\end{hall}

\begin{tamrin}
    What is AX after these instructions? \code{mov ax, 0F0Fh}, then \code{and ax, 00FFh}, then \code{shl ax, 1}.
\end{tamrin}
\begin{hall}
    After \code{and}: \hex{000F}. After the shift (multiply by 2): \hex{001E} = 30.
\end{hall}

\begin{tamrin}
    Write a program that computes $5! = 1 \times 2 \times 3 \times 4 \times 5$ with the \code{loop} instruction, and puts the result in AX.
\end{tamrin}
\begin{hall}
    \asm{ch3/ex_fact.asm}
    The result: AX = 0078 = 120 (tested in the emulator).
\end{hall}

\begin{tamrin}
    After \code{mov sp, 0100h}, then \code{push ax}, then \code{push bx}, then \code{pop cx}: what is SP, and which value is in CX?
\end{tamrin}
\begin{hall}
    After the two pushes SP = \hex{00FC}, and after \code{pop} SP = \hex{00FE}. CX gets the last value in: the value of BX.
\end{hall}

\begin{tamrin}
    What is the vector address of interrupt number \hex{10} (BIOS screen services)?
\end{tamrin}
\begin{hall}
    $\hex{10} \times 4 = 16 \times 4 = 64 = \hex{40}$: the bytes from \hex{00040} to \hex{00043}.
\end{hall}

\begin{tamrin}
    Decode the instruction 89 D8.
\end{tamrin}
\begin{hall}
    89: \code{mov} on 16 bits, d = 0. D8 = \bin{11\,011\,000}: mod = 11 (register), reg = 011 (BX, the source),
    r/m = 000 (AX). So: \code{mov ax, bx}.
\end{hall}
